\documentclass[11pt]{article}
\usepackage[margin=1in]{geometry}
\usepackage{graphicx}
\usepackage{booktabs}
\usepackage{longtable}
\usepackage{amsmath}
\usepackage{hyperref}
\usepackage{float}
\title{Auditor\'ia observacional del Mapper orbital en el cat\'alogo de exoplanetas}
\date{}
\begin{document}
\maketitle

\section{Introducci\'on}
Este subproyecto audita si la topolog\'ia observada en Mapper, en particular en el espacio orbital, est\'a asociada al m\'etodo de descubrimiento de los planetas y no necesariamente a clases f\'isicas intr\'insecas.

\section{Hip\'otesis}
Hip\'otesis A: la estructura topol\'ogica que aparece en Mapper no est\'a determinada principalmente por clases f\'isicas planetarias, sino por sesgos de observaci\'on y m\'etodo de descubrimiento.

\section{Datos y configuraci\'on analizada}
El an\'alisis principal usa la configuraci\'on \texttt{orbital\_pca2\_cubes10\_overlap0p35} ya construida. Cuando hay artifacts disponibles, tambi\'en se resumen \texttt{phys\_min\_pca2\_cubes10\_overlap0p35}, \texttt{joint\_no\_density\_pca2\_cubes10\_overlap0p35}, \texttt{joint\_pca2\_cubes10\_overlap0p35} y \texttt{thermal\_pca2\_cubes10\_overlap0p35}.

\section{Metodolog\'ia}
Para cada nodo $v$ y m\'etodo $m$ se calculan:
\[
p_v(m) = \frac{\mathrm{count}_v(m)}{n_v}
\]
\[
\mathrm{purity}(v) = \max_m p_v(m)
\]
\[
H(v) = -\sum_m p_v(m) \log p_v(m)
\]
\[
H_{norm}(v) = \frac{H(v)}{\log(K)}
\]
donde $K$ es el n\'umero de m\'etodos presentes en el universo analizado.

La prueba nula mantiene fijo el grafo y la membres\'ia nodo-planeta, permuta \texttt{discoverymethod} entre planetas y recalcula m\'etricas globales. El z-score reportado sigue:
\[
z = \frac{\mathrm{observed} - \mathrm{null\_mean}}{\mathrm{null\_std}}
\]

\section{M\'etricas nodo-m\'etodo}
Se generan tablas de pureza, entrop\'ia, m\'etodo dominante, tama\~no nodal, fracci\'on de imputaci\'on media, fracci\'on de variables f\'isicamente derivadas, grado, componente e indicador de periferia.

\section{Prueba nula por permutaci\'on}
Las m\'etricas globales contrastadas contra la distribuci\'on nula son pureza media ponderada, entrop\'ia media ponderada, informaci\'on mutua normalizada entre incidencias nodo-m\'etodo y assortativity del m\'etodo dominante entre nodos conectados cuando existen aristas.

\section{Resultados}
\begin{figure}[H]
\centering
\includegraphics[width=0.9\textwidth]{../../outputs/observational_bias_audit/figures/orbital_graph_by_dominant_discovery_method.pdf}
\caption{Grafo Mapper orbital coloreado por m\'etodo de descubrimiento dominante.}
\end{figure}

\begin{figure}[H]
\centering
\includegraphics[width=0.9\textwidth]{../../outputs/observational_bias_audit/figures/orbital_graph_by_method_purity.pdf}
\caption{Pureza nodal por m\'etodo de descubrimiento.}
\end{figure}

\begin{figure}[H]
\centering
\includegraphics[width=0.9\textwidth]{../../outputs/observational_bias_audit/figures/orbital_graph_by_method_entropy.pdf}
\caption{Entrop\'ia normalizada nodo-m\'etodo en el Mapper orbital.}
\end{figure}

\begin{figure}[H]
\centering
\includegraphics[width=0.9\textwidth]{../../outputs/observational_bias_audit/figures/orbital_node_method_fraction_heatmap.pdf}
\caption{Heatmap nodo x m\'etodo para el espacio orbital.}
\end{figure}

\begin{figure}[H]
\centering
\includegraphics[width=0.75\textwidth]{../../outputs/observational_bias_audit/figures/orbital_permutation_null_nmi.pdf}
\caption{Distribuci\'on nula de NMI y valor observado.}
\end{figure}

\begin{figure}[H]
\centering
\includegraphics[width=0.75\textwidth]{../../outputs/observational_bias_audit/figures/orbital_purity_vs_imputation.pdf}
\caption{Pureza nodal frente a fracci\'on media de imputaci\'on.}
\end{figure}

\begin{figure}[H]
\centering
\includegraphics[width=0.75\textwidth]{../../outputs/observational_bias_audit/figures/orbital_purity_vs_node_size.pdf}
\caption{Pureza nodal frente a tama\~no de nodo.}
\end{figure}

\begin{figure}[H]
\centering
\includegraphics[width=0.9\textwidth]{../../outputs/observational_bias_audit/figures/orbital_component_method_composition.pdf}
\caption{Composici\'on por m\'etodo de descubrimiento a nivel de componente.}
\end{figure}

\begin{figure}[H]
\centering
\includegraphics[width=0.95\textwidth]{../../outputs/observational_bias_audit/figures/config_comparison_bias_metrics.pdf}
\caption{Comparaci\'on resumida de m\'etricas de sesgo observacional entre configuraciones disponibles.}
\end{figure}

\paragraph{Interpretaci\'on autom\'atica.} La NMI observada supera claramente la nula y los nodos con enriquecimiento significativo se concentran en zonas perifericas. La evidencia sugiere que el sesgo observacional existe y parece localizarse sobre todo en ramas o periferias, no necesariamente en el nucleo completo del grafo.

\section{Discusi\'on}
Los resultados deben interpretarse como una auditor\'ia de sesgo observacional sobre grafos ya construidos. No prueban por s\'i mismos la existencia de clases planetarias ni reemplazan una v\'ia contrafactual.

\section{Limitaciones}
Este an\'alisis depende de la calidad de la metadata observacional, del esquema de imputaci\'on ya fijado en el pipeline original y de membres\'ias nodales con solapamiento propio de Mapper. Si nodos de alta pureza tambi\'en tienen alta imputaci\'on, la separaci\'on entre sesgo observacional y dependencia de imputaci\'on requiere an\'alisis adicional.

\section{Conclusi\'on}
El subproyecto deja cuantificada la asociaci\'on entre topolog\'ia Mapper y m\'etodo de descubrimiento en t\'erminos de pureza, entrop\'ia, enriquecimiento local y contraste por permutaciones.

\appendix
\section{Ap\'endice de tablas}
\input{../../outputs/observational_bias_audit/tables/summary_global_bias_metrics.tex}
\input{../../outputs/observational_bias_audit/tables/top_enriched_nodes.tex}
\input{../../outputs/observational_bias_audit/tables/central_vs_peripheral_bias.tex}

\end{document}
